% page 246 of the facsimile (image p-245 of the scan)
% block 170.2 x 242.0 mm ; left margin 8.0 mm
\pgc{-12.700mm}{0.000mm}{
\l{238}
\l{}
\l{\sou{instanto}.\cel{2}Mikra tempo-fragmento. - DEFIRS.}
\l{}
\l{\sou{instigar}. (\sou{trans., ad)}.\cel{2}Pulsar ad ago. - EFIS.}
\l{}
\l{\sou{instilar}. (\sou{trans. aden}) (\sou{medic.})\cel{2}Varsar gutope. - DEFI.}
\l{}
\l{\sou{instinto}. Impulsivo naturala ; impulsivo interna qua insti-}
\l{\cel{3}gas ento vivanta (homo o bestio) ad agi nerezonita, por}
\l{\cel{3}la konservo di la individuo, di la speco. - DEFIRS.}
\l{}
\l{\sou{institucar}. (\sou{trans.}) Establisar durive. - DEFIRS.}
\l{}
\l{\sou{instituto}. Korpo institucura ek literaturisti, ciencisti,}
\l{\cel{3}artisti, e c. - DEFIRS.}
\l{}
\l{\sou{instruciono}. Expliko por la direkto di afero, por la uzo di}
\l{\cel{3}ulo. - DEFIRS.}
\l{}
\l{\sou{instruktar}. (\sou{trans.}). Formacar la spirito,la mento, la psiko}
\l{\cel{3}di (ulu) per lecioni, precepti. - DEFIRS.}
\l{}
\l{\sou{instrumento}.\cel{2}To quon onu uzas por produktar efekto materiala.}
\l{\cel{3}- DEFIRS.}
\l{}
\l{\sou{insulo}.\cel{2}Extensajo ek tero qua,en maro, cirkondesas da aquo}
\l{\cel{3}omna-latere, tota-cirkume. - DEFIS.}
\l{}
\l{\sou{insultar}. (\sou{trans}.) Ofensar extreme (generale per vorti).DEFIS.}
\l{}
\l{\sou{-int-}. (\sou{gram.})\cel{2}Dezinenco di la participo pasintala aktiva.}
\l{}
\l{\sou{intalio}. Petro harda, grabita kave. - DEFI.}
\l{}
\l{\sou{integra}. Di qua la toto ne subisis diminuto. - EFIRS.}
\l{}
\l{\sou{integralo.}\cel{2}(\sou{mat}.)La integralo di diferencialo :quanto di qua}
\l{\cel{3}la diferencialo esas infinitezima. - DEFIRS.}
\l{}
\l{\sou{intelektar. (trans.})\cel{2}Konoceskar per la koncepto-fakultato}
\l{\cel{3}pura. - DEFIRS.}
\l{}
\l{\sou{intelektuelo.}\cel{3}Persono qua prizas prepondere lo pensala, lo}
\l{\cel{3}spiritala. - DEFIS.}
\l{}
\l{\sou{inteligenta}. Qua komprenas facile. - DEFIRS.}
\l{}
\l{\sou{intencar}. (\sou{trans.)}\cel{2}Volar ad skopo. - FIS.}
\l{}
\l{\sou{intendanto}.\cel{2}(\sou{armeo).}\cel{2}Ensemblo ek la oficiri qui su okupas pri}
\l{\cel{3}la\cel{2}servicii administrala di la armearo. - DFIRS.}
\l{}
\l{\sou{intendanto}. I. La persono di qua la tasko konsistas ye jerar}
\l{\cel{3}la aferi, administrar la menajo di richo, di sinioro. - II.}
\l{\cel{4}Stat-oficiero, chefo di administrantaro guvernantarala. DEIFRS.}
}
